\section{Related Work}
\label{sec:related_work}

\begin{table*}[!t]
\caption{Taxonomy of Representative Methodologies for VRPTW and Related Routing Problems}
\label{tab:taxonomy}
\centering
\footnotesize
\setlength{\tabcolsep}{3.2pt}

\begin{tabularx}{\textwidth}{
    l
    X
    c
    c
    c
    c
    c
}
\toprule

\textbf{Methodology Category}
&
\textbf{Representative Works}
&
\shortstack{\textbf{Hard TW}\\\textbf{Feasible}}
&
\shortstack{\textbf{Objective}\\\textbf{Paradigm}}
&
\shortstack{\textbf{State-Aware}\\\textbf{Selection}}
&
\shortstack{\textbf{Hierarchical}\\\textbf{Control}}
&
\shortstack{\textbf{Learned}\\\textbf{Accept}}
\\

\midrule

\multirow{2}{*}{\textbf{Classical Metaheuristics}}
&
ALNS \cite{Ropke2006, pisinger2007general}
&
$\checkmark$
&
Lexicographic
&
$\times$ (Stateless)
&
$\times$
&
$\times$ (SA)
\\

&
SISR \cite{christiaens2020slack}, HGS \cite{vidal2013hybrid}
&
$\checkmark$
&
Lexicographic
&
$\times$ (Rule-based)
&
$\times$
&
$\times$ (SA/Threshold)
\\

\addlinespace[3pt]

\multirow{2}{*}{\textbf{End-to-End Neural Solvers}}
&
AM \cite{kool2019attention}, DRL-VRP \cite{nazari2018reinforcement}
&
$\times$ (Penalized)
&
Scalarized
&
N/A
&
$\times$
&
N/A
\\

&
Neural VRPTW \cite{falkner2020learning, lin2021deep}
&
$\times$ (Mask-fragile)
&
Scalarized
&
N/A
&
$\times$
&
N/A
\\

\addlinespace[3pt]

\multirow{2}{*}{\textbf{Single-Agent Hybrid RL-LNS}}
&
Learned LNS \cite{lu2020learning, son2023learning}
&
$\checkmark$
&
Scalarized
&
$\checkmark$
&
$\times$ (Single-agent)
&
$\times$ (SA)
\\

&
Neural LNS \cite{hottung2020neural}
&
$\checkmark$
&
Scalarized
&
$\checkmark$
&
$\times$ (Single-agent)
&
$\times$ (Greedy)
\\

\addlinespace[3pt]

\midrule

\textbf{Proposed Hybrid DDQN-ALNS}
&
\textbf{This Work}
&
$\mathbf{\checkmark}$
&
\textbf{Lexicographic}
&
$\mathbf{\checkmark}$
&
$\mathbf{\checkmark}$ \textbf{(Dual HMDP)}
&
$\mathbf{\checkmark}$ \textbf{(LAC)}
\\

\bottomrule
\end{tabularx}
\end{table*}

\subsection{Classical Metaheuristics and ALNS for VRPTW}
\label{sec:rw_classical}

Metaheuristics have long dominated the practical resolution of the VRPTW
\cite{braysy2005vehicle}. Building on the seminal insertion heuristics of
Solomon \cite{Solomon1987} and the Large Neighborhood Search (LNS)
paradigm of Shaw \cite{shaw1998using}, Ropke and Pisinger
\cite{Ropke2006, pisinger2007general, pisinger2010large} formalized the
Adaptive Large Neighborhood Search (ALNS) framework, in which a portfolio of
handcrafted removal and insertion operators competes for selection via
roulette-wheel probabilities updated from historical performance scores.

Subsequent research advanced two complementary directions. For the
hierarchical objective of fleet minimization, Guided Ejection Search
\cite{nagata2009ges} and penalty-based edge assembly memetic algorithms
\cite{nagata2010srex} remain the strongest dedicated mechanisms. In parallel,
matheuristic hybrids exploit an elite \emph{route pool}---tracing back to the
adaptive memory of Rochat and Taillard \cite{rochat1995probabilistic}---and
periodically solve a Set Partitioning model over accumulated routes
\cite{alvarenga2007genetic, subramanian2013hybrid, vidal2013hybrid},
recombining high-quality route fragments beyond the reach of local moves.
More recently, Slack Induction by String Removals \cite{christiaens2020slack}
demonstrated that carefully designed destroy operators alone can rival far
more complex frameworks, underscoring how decisive the destroy step is.

Despite this empirical success, the adaptive layer of ALNS exhibits an
intrinsic architectural limitation. Operator weights are updated as
exponentially smoothed averages over coarse segments of $\eta$ iterations,
yielding credit assignment that is both delayed and aggregated; critically,
the resulting selection distribution is \emph{stateless}, conditioning on
neither the spatial topology of the incumbent (e.g., cluster structure,
inter-route proximity) nor its temporal characteristics (e.g., time-window
tightness, slack distribution). Performance is further sensitive to manually
tuned hyperparameters such as the decay factor and reward scores. Notably, a
recent meta-analysis by Turke\v{s} et al.\ \cite{turkes2021meta} quantified
that this adaptive mechanism contributes only marginal gains over
non-adaptive LNS---strong evidence that the operator-selection policy, rather
than the operator set itself, is the underexploited component of the
framework.

\subsection{Neural Combinatorial Optimization and End-to-End Deep RL}
\label{sec:rw_nco}

In parallel with classical heuristics, Neural Combinatorial Optimization (NCO)
has gained widespread traction by learning data-driven routing policies.
Originating from Pointer Networks \cite{vinyals2015pointer} and policy-gradient
formulations for the Traveling Salesperson Problem \cite{bello2016neural},
Nazari et al.\ \cite{nazari2018reinforcement} extended end-to-end reinforcement
learning to the Capacitated VRP by decoupling static customer coordinates from
dynamic demand states. Subsequently, Kool et al.\ \cite{kool2019attention}
established the Attention Model (AM) as the prevailing benchmark, leveraging
Transformer-style self-attention with a greedy rollout baseline. Graph Neural
Networks (GNNs) \cite{khalil2017learning, joshi2019efficient} have likewise been
deployed to capture permutation-invariant spatial graph embeddings.

More recently, several architectures have attempted direct end-to-end
construction for VRPTW \cite{falkner2020learning, lin2021deep}. Despite their
impressive inference speed on small-scale benchmark instances, pure
end-to-end neural solvers face four fundamental bottlenecks that impede their
deployment in industrial VRPTW settings:

\begin{enumerate}
    \item \textbf{Feasibility Fragility under Tight Time Windows}: While capacity
    constraints can be masked during autoregressive decoding, hard time windows
    frequently induce structural deadlocks where no remaining unvisited customer
    is temporally reachable, forcing the constructive policy to spawn redundant
    vehicles or incur severe penalty violations.
    \item \textbf{Lexicographic Misalignment}: Pure DRL frameworks are trained
    to minimize a scalarized objective (typically travel distance or a
    weighted surrogate), lacking the mechanism to strictly enforce the
    lexicographic priority ($\NV \succ \TD$) that governs fleet economics.
    \item \textbf{Out-of-Distribution (OOD) Scalability Collapse}: Models trained
    on synthetic uniform distributions at $N \in \{50, 100\}$ exhibit steep
    performance degradation when evaluated on larger problem instances
    ($N \in \{200, 400\}$) or heterogeneous spatial clusters without expensive
    instance-specific fine-tuning.
    \item \textbf{Lack of Operational Auditability}: Constructing an entire routing
    plan via end-to-end neural forward passes functions as an opaque black box,
    offering no interpretable intermediate steps for dispatchers to audit,
    constrain, or manually adjust in real-time dispatch systems.
\end{enumerate}

\subsection{Learning-Augmented Metaheuristics}
\label{sec:rw_hybrid}

To circumvent the feasibility and scalability bottlenecks of pure neural construction, the ``Learning to Optimize'' paradigm \cite{bengio2021machine} embeds machine learning models directly within proven metaheuristic frameworks. Rather than generating solutions from scratch, learning agents guide critical decision points inside iterative search loops, such as operator selection \cite{syberfeldt2009reinforcement, lu2020learning}, neural solution perturbation \cite{chen2019learning}, and learned destroy policies \cite{hottung2020neural, son2023learning}.

While these hybrid frameworks preserve the underlying validity of metaheuristic search, existing learning-augmented approaches for routing exhibit three key shortcomings:
\begin{enumerate}
    \item \textbf{Single-Level Decision Architecture}: Prior RL-augmented LNS methods deploy a flat, single-agent policy that selects operators uniformly across all search stages. They lack hierarchical decoupling between macro-level exploration regimes (e.g., detecting search stagnation and triggering aggressive plateau escape) and micro-level operator selection conditioned on local topology.
    \item \textbf{Persistence of Static Acceptance Criteria}: Despite deploying neural operator selectors, existing hybrids continue to rely on conventional Simulated Annealing with static geometric cooling schedules. Consequently, candidate moves are judged solely by instantaneous cost degradation $\Delta C$, ignoring structural potential that could unlock downstream improvements.
    \item \textbf{Absence of Lexicographic Multi-Objective Coordination}: Most learned metaheuristics focus on scalarized objectives (e.g., distance minimization in CVRP), without architectural mechanisms designed for the strict lexicographic priority ($\NV \succ \TD$) fundamental to VRPTW fleet operations.
\end{enumerate}

\subsection{Positioning of This Work}
\label{sec:rw_positioning}

Table~\ref{tab:taxonomy} situates the proposed Hybrid DDQN-ALNS framework within the landscape of vehicle routing methodologies. Our approach distinctively combines the hard constraint guarantees of classical operations research with a dual-agent Hierarchical MDP and a Learned Acceptance Criterion, providing a rigorous, state-aware, and lexicographically aligned optimization engine.
